%& -shell-escape
% !TeX root = thesis-text.tex
% !TeX encoding = UTF-8
% !TeX program = XeLaTeX

\documentclass{.local/lib/classes/ndvr-super-article-standalone}

\title[NDVR thesis-abstract]{%
    Проблема поиска нечетких дубликатов видео%
}

\begin{document}

    \section{Системы рекомендаций}

    Многие сервисы в~Интернете используют рекомендательные системы,
    которые пытаются предсказать, какие объекты (музыка, книги, новости)
    будут интересны пользователю.
    Если зрителю понравился фильм,
    он~будет интересоваться и~похожими фильмами \cite{Cherubini:2009},
    т.~е. нечеткими дубликатами этого фильма.
    Используя эту особенность, библиотеки видео
    могут более точно подстраиваться под предпочтения зрителя.

    Тем не~менее, текущие системы рекомендации,
    как~правило, используют только текстовую информацию:
    названия, комментарии, описания и~другие метаданные.
    Т.~к. абсолютно разные видео могут иметь похожие описания,
    то текстовый подход вносит значительный шум. 
    Пример неудачной рекомендации сервиса YouTube показан
    на рисунке \ref{fig:video-recommendation}.
    На основе просмотра сюжета про падение Челябинского метеорита\footnote{
        Челябинский метеорит~— космическое тело, 
        которое упало на земную поверхность 
        15 февраля 2013 года.
        Наиболее крупные фрагменты метеорита были найдены 
        на дне озера Чебаркуль в Челябинской области.
    } зрителю были предложены не релевантные 
    сюжеты про электричку «Миасс-Челябинск»
    и про рыбную ловлю на озере Чебаркуль.
    
    \begin{figure}[H]
        \includegraphics[width=0.7\textwidth]%
            {.local/img/video/apps/video-recommendation.png}
        \caption{Пример ошибочной рекомендации видео}
        \label{fig:video-recommendation}
    \end{figure}

    Анализ содержимого видео 
    и применение методов поиска нечётких дубликатов 
    в~таких системах может увеличить точность рекомендаций.
    Для~этого достаточно сопоставлять текстовую 
    информацию с~самим видео \cite{Cherubini:2009}.
    При этом предполагается, что некоторых исходный набор 
    видео должен быть заранее размечен.
    Смысловую информацию в видео могут добавлять
    пользователи сервиса. В работе \cite{Zhao:2010} предлагают 
    мотивировать пользователей с помощью элементов социальных сетей.

\end{document}
